\section*{Chapitre \ref{ch:b1:proteins} --- Acides aminés et protéines}

\begin{solution}{exo:b1:proteins:1}
$\mathrm{^+H_3N{-}CH(R){-}COO^-}$~: un \omterm{def:b1:proteins:aminoacid}{zwitterion}. Leu apolaire
aliphatique~; Ser \omterm{def:b1:water-small-molecules:water}{polaire} non chargée~; Lys chargée positivement~; Glu
chargée négativement~; Phe aromatique.
\end{solution}

\begin{solution}{exo:b1:proteins:2}
Primaire~: la séquence, \omterm{def:b1:proteins:peptide}{liaisons peptidiques}. Secondaire~: le repliement
local du squelette (hélice, feuillet), \omterm{def:b1:water-small-molecules:water}{liaisons hydrogène} du squelette.
Tertiaire~: le repliement de la chaîne entière~; \omterm{prop:b1:water-small-molecules:properties}{effet hydrophobe},
\omterm{def:b1:water-small-molecules:water}{liaisons hydrogène}, ponts salins, contacts de \omterm{def:b1:water-small-molecules:bonds}{van der Waals}, \omterm{def:b1:proteins:tertiary}{ponts
disulfure}. Quaternaire~: l'assemblage des sous-unités, par les mêmes
forces non covalentes.
\end{solution}

\begin{solution}{exo:b1:proteins:3}
$450\times 110 = \qty{50}{kDa}$. Hélice~: $450\times 0.15 = \qty{68}{nm}$.
Brin~: $450\times 0.35 = \qty{158}{nm}$.
\end{solution}

\begin{solution}{exo:b1:proteins:4}
Environ $0.98$ et $0.76$~: un cinquième de la charge (\qty{22}{\%}) est
délivré.
\end{solution}

\begin{solution}{exo:b1:proteins:5}
Groupes~: aminé N-terminal (9.5), \omterm{def:b1:proteins:aminoacid}{chaîne latérale} de Lys (10.5), \omterm{def:b1:proteins:aminoacid}{chaîne
latérale} d'Asp (3.9), carboxyle C-terminal (2). pH 7~:
$+1 + 1 - 1 - 1 = 0$. pH 2~: $+1 + 1 + 0 - 0.5 = +1.5$ (le carboxyle
terminal à moitié ionisé). pH 12~: $0 + 0 - 1 - 1 = -2$ (Lys en grande
partie déprotonée). Charge nulle entre le p$K_a$ d'Asp (3.9) et celui du
groupe aminé (9.5), où la charge vaut $0$ tout du long~:
p$I \approx 6.7$ (moyenne de 3.9 et de 9.5).
\end{solution}

\begin{solution}{exo:b1:proteins:6}
La \omterm{def:b1:proteins:aminoacid}{chaîne latérale} de la proline est refermée sur son propre azote~:
l'azote ne porte plus d'hydrogène à donner à la \omterm{def:b1:water-small-molecules:water}{liaison hydrogène} de
l'hélice, et $\phi$ est bloqué vers $-60^\circ$, ce qui n'est toléré
qu'au premier tour d'une hélice. La glycine, dont la \omterm{def:b1:proteins:aminoacid}{chaîne latérale} est
un hydrogène, ne subit aucune gêne stérique et peut adopter des valeurs
de $\phi$ et de $\psi$ interdites à tout autre résidu, dont celles des
coudes serrés.
\end{solution}

\begin{solution}{exo:b1:proteins:7}
Appariements de huit cystéines~: $7\times 5\times 3\times 1 = 105$~; une
réoxydation au hasard donne le jeu natif environ \qty{1}{\%} du temps.
La récupération d'une activité presque complète a montré que c'était la
chaîne, et non le hasard, qui choisissait l'appariement~: la séquence
code le repliement et les \omterm{def:b1:proteins:tertiary}{ponts disulfure} ne font que le verrouiller.
\end{solution}

\begin{solution}{exo:b1:proteins:8}
$P_{50}^{2.8} = 33.4$. À 2~: $2^{2.8} = 6.96$, $Y = 0.17$~; à 3.5~:
$0.50$~; à 8~: $8^{2.8} = 337$, $Y = 0.91$. Avec $n = 1$~: $Y =
P/(3.5 + P)$~: $0.36$, $0.50$, $0.70$. Livraison entre 8 et 2~:
coopératif $0.74$, non coopératif $0.34$~: la \omterm{def:b1:proteins:peptide}{protéine} \omterm{def:b1:proteins:allostery}{coopérative}
délivre deux fois plus.
\end{solution}

\begin{solution}{exo:b1:proteins:9}
Un tétramère de quatre sous-unités identiques de \qty{50}{kDa}. Faire
migrer le gel SDS sans réducteur~: si des bandes apparaissent à 100, 150
ou \qty{200}{kDa}, des \omterm{def:b1:proteins:tertiary}{ponts disulfure} relient les sous-unités~; s'il ne
reste que la bande de \qty{50}{kDa}, elles sont tenues de façon non
covalente.
\end{solution}

\begin{solution}{exo:b1:proteins:10}
Le glutamate est chargé et hydraté à la surface~; la valine est
hydrophobe et crée une plage collante que l'eau exclut. Dans l'état T
désoxygéné, une poche hydrophobe complémentaire est exposée sur une
molécule voisine, de sorte que les molécules se collent bout à bout en
fibres qui déforment le globule rouge. Dans les poumons, l'état R cache
la poche et les fibres se dissolvent~; dans les \omterm{def:b1:body-plans-tissues:tissue}{tissus}, où l'hémoglobine
est désoxygénée, elles se forment --- et c'est là que les globules se
falciforment et bloquent les petits vaisseaux.
\end{solution}

\begin{solution}{exo:b1:proteins:11}
À \qty{37}{\celsius}, $RT = \qty{2.58}{kJ/mol}$~: $K = e^{-40/2.58} =
e^{-15.5} = \num{1.8e-7}$~: une molécule sur cinq millions est dépliée.
À \qty{45}{\celsius} avec \qty{10}{kJ/mol}, $RT = 2.64$~: $K = e^{-3.8} =
0.022$~: \qty{2}{\%} de dépliées, et la proportion monte vite. Une
stabilité marginale permet aux \omterm{def:b1:proteins:peptide}{protéines} de changer de forme quand elles
travaillent (\omterm{def:b1:proteins:allostery}{allostérie}, catalyse), d'être dépliées pour traverser des
membranes, et d'être dégradées et remplacées vite quand elles sont
endommagées ou devenues inutiles.
\end{solution}

\begin{solution}{exo:b1:proteins:12}
Anfinsen a montré que la séquence suffit à spécifier le repliement, et
l'\omterm{def:b1:proteins:allostery}{allostérie} montre que la forme --- et ses changements --- est la
fonction. Mais la \omterm{def:b1:cell-unit-of-life:cell}{cellule} intervient à chaque étape~: les \omterm{prop:b1:proteins:denaturation}{chaperons}
sauvent les chaînes qui s'agrégeraient avant de se replier, car le
\omterm{def:b1:eukaryotic-cell:organelle}{cytosol} encombré n'est pas un tube à essai dilué~; une seule
substitution (drépanocytose) donne une \omterm{def:b1:proteins:peptide}{protéine} correctement repliée
dont la surface est fautive, de sorte que la «~forme~» doit inclure la
chimie de surface et non le seul squelette~; et beaucoup de \omterm{def:b1:proteins:peptide}{protéines}
exigent des modifications, des partenaires ou des ligands ajoutés après
le repliement avant de fonctionner. La phrase énonce le principe~; la
\omterm{def:b1:cell-unit-of-life:cell}{cellule} en fournit les conditions.
\end{solution}

\begin{solution}{pb:b1:proteins:1}
\textbf{1.} $150\times 1.34 = \qty{201}{mL}$ d'$\mathrm{O_2}$ par litre.
\textbf{2.} $e^{3.51} = 33.4$~; $e^{4.67} = 106.7$~; $e^{7.25} = 1408$.
\textbf{3.} Poumons~: $1408/(33.4 + 1408) = 0.977$. \omterm{def:b1:body-plans-tissues:tissue}{Tissus}~:
$106.7/(33.4 + 106.7) = 0.762$.
\textbf{4.} $0.977 - 0.762 = 0.215$~: $0.215\times 201 = \qty{43}{mL}$
par litre.
\textbf{5.} $250/43 = \qty{5.8}{L/min}$ --- proche du débit cardiaque
de repos.
\textbf{6.} Myoglobine~: poumons $13.3/13.67 = 0.973$~; \omterm{def:b1:body-plans-tissues:tissue}{tissus}
$5.3/5.67 = 0.935$~; elle délivre $0.038\times 201 = \qty{7.7}{mL}$ par
litre, six fois moins~: un transporteur hyperbolique qui se charge bien
ne sait pas décharger aux pressions des \omterm{def:b1:body-plans-tissues:tissue}{tissus}.
\textbf{7.} La \omterm{def:b1:proteins:allostery}{coopérativité} place la partie raide de la courbe entre
les deux pressions de travail, de sorte qu'une faible baisse de pression
libère une grande fraction de la charge.
\textbf{8.} $e^{2.78} = 16.1$~; $Y = 16.1/(33.4 + 16.1) = 0.325$.
\textbf{9.} $e^{4.21} = 67.4$~; $Y = 16.1/(67.4 + 16.1) = 0.193$.
\textbf{10.} Sans Bohr~: $(0.977 - 0.325)\times 201 = \qty{131}{mL}$~;
avec~: $(0.977 - 0.193)\times 201 = \qty{158}{mL}$ par litre~; gain
\qty{27}{mL}, soit \qty{20}{\%}.
\textbf{11.} $2.7/3.07 = 0.88$ de saturation~; à \qty{0.5}{kPa},
$0.5/0.87 = 0.57$~: elle libère le tiers de sa réserve vers les
\omterm{def:b1:eukaryotic-cell:mitochondrion}{mitochondries} pendant la contraction, et se recharge entre les
contractions.
\textbf{12.} Les protons se fixent sur des sites (histidines, extrémités
N-terminales) éloignés des hèmes et, ce faisant, stabilisent l'état T,
ce qui abaisse l'affinité des hèmes~: un ligand sur un site change
l'affinité en un autre --- c'est l'\omterm{def:b1:proteins:allostery}{allostérie}.
\textbf{13.} $e^{5.45} = 233$~; poumons $233/(33.4 + 233) = 0.875$~;
\omterm{def:b1:body-plans-tissues:tissue}{tissus} $0.762$~; livraison $0.113\times 201 = \qty{23}{mL}$~: environ la
moitié du niveau de la mer.
\textbf{14.} $e^{3.88} = 48.4$~: poumons $233/(48.4 + 233) = 0.828$~;
\omterm{def:b1:body-plans-tissues:tissue}{tissus} $106.7/(48.4 + 106.7) = 0.688$~; livraison $0.140\times 201 =
\qty{28}{mL}$.
\textbf{15.} La charge baisse un peu (de 0.875 à 0.828) mais la décharge
baisse davantage (de 0.762 à 0.688), de sorte que la différence
augmente~: à \qty{7}{kPa} les poumons sont encore sur le plateau de la
courbe alors que les \omterm{def:b1:body-plans-tissues:tissue}{tissus} sont sur sa partie raide.
\textbf{16.} Capacité $180\times 1.34 = \qty{241}{mL}$~; $0.140\times
241 = \qty{34}{mL}$ par litre --- les trois quarts du niveau de la mer
retrouvés.
\textbf{17.} $6^{2.8} = e^{2.8\times 1.79} = e^{5.02} = 151$~: poumons
$233/384 = 0.607$, \omterm{def:b1:body-plans-tissues:tissue}{tissus} $106.7/258 = 0.414$~: livraison $0.193$,
meilleure encore sur le papier --- mais le sang artériel ne serait
saturé qu'à \qty{61}{\%}, et toute baisse supplémentaire de la pression
pulmonaire ou toute exigence d'une pression tissulaire plus élevée
laisserait l'encéphale à court~; l'ajustement réel s'arrête vers
\qty{4}{kPa}.
\textbf{18.} $48.4/(33.4 + 48.4) = 0.59$.
\textbf{19.} $e^{2.68} = 14.6$~; $48.4/(14.6 + 48.4) = 0.77$.
\textbf{20.} À pression égale, la \omterm{def:b1:proteins:peptide}{protéine} fœtale fixe davantage~: à
mesure que le sang de la mère décharge vers \qty{59}{\%} et que le sang
fœtal se charge vers \qty{77}{\%}, l'oxygène descend le gradient de
pression que la différence d'affinités entretient de part et d'autre de
la membrane placentaire.
\textbf{21.} Le BPG se lie à l'état T et abaisse l'affinité~; les
chaînes $\gamma$ le lient faiblement, de sorte que l'hémoglobine fœtale
reste plus près de sa forte affinité intrinsèque~: $P_{50}$ plus bas.
\textbf{22.} Poumons~: $1408/(14.6 + 1408) = 0.990$~; \omterm{def:b1:body-plans-tissues:tissue}{tissus} $106.7/(14.6
+ 106.7) = 0.880$~; livraison $0.110\times 201 = \qty{22}{mL}$, la
moitié de celle de l'adulte~: bon pour prendre l'oxygène à la mère,
mauvais pour le délivrer une fois les poumons en service, d'où le
changement.
\textbf{23.} $5^{2.8} = e^{2.8\times 1.609} = e^{4.51} = 90.6$~: poumons
$1408/1499 = 0.94$, \omterm{def:b1:body-plans-tissues:tissue}{tissus} $106.7/197 = 0.54$~: livraison $0.40$, presque
le double de la normale --- la personne est légèrement cyanosée aux
lèvres mais le tolère bien~; le prix est une réserve moindre en
altitude.
\textbf{24.} $Y = P/(3.5 + P)$~: poumons $0.79$, \omterm{def:b1:body-plans-tissues:tissue}{tissus} $0.60$~:
livraison $0.19\times 201 = \qty{38}{mL}$ --- et, dans le muscle qui
sprinte à \qty{2.7}{kPa}, $0.44$~: livraison $0.35\times 201 = \qty{70}{mL}$
contre \qty{158}{mL}. Sans \omterm{def:b1:proteins:allostery}{coopérativité}, le transporteur se charge mal
et décharge mal.
\textbf{25.} \qty{43}{mL} d'$\mathrm{O_2}$ par litre au repos,
\qty{158}{mL} par litre dans un muscle qui sprinte~; la différence vient
de la \omterm{def:b1:proteins:allostery}{coopérativité} (la courbe sigmoïde) et de l'effet Bohr (le décalage
de $P_{50}$ en milieu acide).
\end{solution}
